\honest{On Expressing Your Concerns}{Eliezer Yudkowsky}{2007}

Asch's experiments have a scary side, a hopeful side and a wearisome side. A room of people saying black is white makes many others say it too; a single dissenter drives conformity down; and dissent ``was not learned,'' because when the dissenter sided with the group, conformity rose again. \nb{Asch put that rise down to the partner's ``desertion.'' When the partner simply left the room, ``the partner's effect outlasted his presence.''} Dissent helps the group, costs the dissenter, has to be kept up, and may be wrong.

Recently I told two people planning a project that they were too optimistic. A fourth man, who had come to give one of them a ride, answered every problem I raised with ``Don't worry, I'm sure we can handle it!'' From this I had a sudden insight. Alone, people doubt themselves. Together, politeness forbids doubting each other, so each takes the other's reassurance for confidence. In a whole group this is ``pluralistic ignorance,'' which I call the most fearsome possibility raised by Asch's experiments.

Robin Hanson and I disagree about when to disagree, but we both hold that Aumann's theorem ``extends to imply that common knowledge of a factual disagreement shows someone must be irrational.'' \nb{The theorem assumes the parties start from the same priors, as Hanson's own statements of the result say. The step to ``someone must be irrational'' needs the further claim that rational agents must share priors, which Hanson has argued and others dispute. The post mentions neither the condition nor the dispute.} So for rationalists, disagreeing is serious business.

The most important lesson of Asch's experiments, I think, is to distinguish expressing a concern from disagreeing. Ideal Bayesians converge by sharing evidence the listener could not predict. But socially there is not much difference between the two. Speaking out cannot be undone, and a group of rationalists could only ``agree to pretend'' there is a difference. Cynical self-help books, such as Machiavelli's \textsc{The Prince}, advise you to mask your nonconformity entirely. Do not expect thanks.

In the end, ``the decision is up to you.'' I give no advice on when a concern is worth its cost.
